\begin{figure}[ht]
\centering
\begin{tikzpicture}[>=Stealth, line width=.7pt,
  every node/.style={font=\small},
  endpoint/.style={circle,fill=black,inner sep=1.5pt},
  pathlabel/.style={fill=white,inner sep=2pt}]
  \begin{scope}
    \node at (2.45,2.15) {Before};
    \coordinate (a) at (0,0);
    \coordinate (b) at (1.35,0);
    \coordinate (c) at (3.55,0);
    \coordinate (d) at (4.90,0);
    \draw[->] (a) -- node[pathlabel,above] {$Q_1$} (b);
    \draw[->,very thick] (b) -- node[pathlabel,above] {$I$} (c);
    \draw[->] (c) -- node[pathlabel,above] {$Q_2$} (d);
    \foreach \q in {a,b,c,d} \node[endpoint] at (\q) {};
    \node[below=5pt] at (a) {$q_-$};
    \node[below=5pt] at (b) {$i_-$};
    \node[below=5pt] at (c) {$i_+$};
    \node[below=5pt] at (d) {$q_+$};
    \node at (2.45,-.88) {$Q=Q_1IQ_2,\qquad |E(I)|\ge31$};
  \end{scope}
  \begin{scope}[xshift=6.2cm]
    \node at (2.45,2.15) {After};
    \coordinate (a) at (0,0);
    \coordinate (b) at (1.35,0);
    \coordinate (c) at (3.55,0);
    \coordinate (d) at (4.90,0);
    \draw[->] (a) -- node[pathlabel,above] {$Q_1$} (b);
    \draw[->] (c) -- node[pathlabel,above] {$Q_2$} (d);
    \draw[->,very thick] (a) to[out=65,in=115]
      node[pathlabel,above] {$Z,\quad |E(Z)|\le12$} (d);
    \foreach \q in {a,b,c,d} \node[endpoint] at (\q) {};
    \node[below=5pt] at (a) {$q_-$};
    \node[below=5pt] at (b) {$i_-$};
    \node[below=5pt] at (c) {$i_+$};
    \node[below=5pt] at (d) {$q_+$};
    \node at (2.45,-.88) {Bypass from $i_-$ to $i_+$: $Q_1^{-1}ZQ_2^{-1}$};
  \end{scope}
\end{tikzpicture}
\caption{The shortening move: add $Z$ between the endpoints of $Q$,
then delete the edges and interior vertices of $I$. The retained paths
$Q_1,Q_2$ and the new path give a bypass with the same value in $D$
as $I$. The drawing is schematic: displayed endpoints may coincide,
and $Q_1,Q_2$ may overlap or be empty. The proof establishes that both
avoid the edges and interior vertices of $I$. Each line denotes a path,
and the rest of the graph is omitted.}
\label{rank:shortening-schematic}
\end{figure}
